% MoRE architecture paper -- WORKING DRAFT
%
% Status: skeleton with committed framing. Prose is real where the claim is
% already decided; every number is a placeholder until the ablation campaign
% reports (see ablation_campaign.md).
%
% TONE POLICY (deliberate -- keep it consistent as the draft grows):
%   Humour is allowed in the abstract (one dry line, max), the introduction,
%   figure captions, and the limitations section. Methods, results, and related
%   work stay sober -- a joke in a results table reads as a hedge.
%   One hard rule: never joke at another paper's expense. Joke at the problem,
%   at uniform compute, and at ourselves. The authors of the work we position
%   against are plausibly our reviewers.
%
% Conventions:
%   depth = passes = recursions + 1   (depth 1 is the non-looped base case)
%   k     = routed experts per pass; one shared expert is always active
%   \TODO{...} marks something that must be written or measured before submission
%   \CITE{...} marks a citation that must be resolved
%
% Title candidates (swap freely, decided before first public artifact):
%   [chosen] Three Dials, One Budget: Per-Token Compute Allocation with MoRE
%   MoRE Is More: Joint Per-Token Routing of Depth, Width, and Pathway
%   How Much Compute Does a Token Deserve?
%   Goldilocks Routing: Never Too Much, Never Too Little, Per Token
%
\documentclass[11pt]{article}

\usepackage[margin=1in]{geometry}
\usepackage{array}
\usepackage{booktabs}
\usepackage{longtable}
\usepackage{hyperref}
\usepackage{amsmath}
\usepackage{amssymb}
\usepackage{xcolor}

\newcommand{\code}[1]{\texttt{\detokenize{#1}}}
\newcommand{\TODO}[1]{\textcolor{red}{[TODO: #1]}}
\newcommand{\CITE}[1]{\textcolor{blue}{[cite: \detokenize{#1}]}}
\newcommand{\kmax}{k_{\max}}
\newcommand{\Rmax}{R_{\max}}
\setlength{\parskip}{0.55em}
\setlength{\parindent}{0pt}
\sloppy

\title{Three Dials, One Budget:\\Per-Token Compute Allocation with MoRE}
\author{Giulianno Vollmer\\Lernex\\\texttt{giulianno@lernex.net}}
\date{\TODO{date}}

\begin{document}
\maketitle

\begin{abstract}
Language models spend the same computation on the word \emph{the} as on the step
of a proof where everything goes wrong. Conditional-compute architectures fix
this, but each of them fixes it along exactly one axis: Mixture-of-Depths and
Mixture-of-Recursions vary \emph{how many} steps a token gets, adaptive-$k$
Mixture-of-Experts designs vary \emph{how wide} a single step is, and recent
looped MoE models fix both and vary neither.
We present \textbf{MoRE}, in which a single per-token compute budget is allocated
jointly across three axes: recursion \textbf{depth} (how many passes through a
weight-shared stack), expert \textbf{width} (how many experts at each pass), and
expert \textbf{pathway} (which experts, re-decided every pass from the token's
evolving state). All three decisions are learned end to end from task loss under
shared budget objectives. There is no halting threshold to tune, because there is
no halting threshold.
We add \textbf{MoRE-RM}, a recurrent depth memory whose entries are typed by
pass, anchor, expert coalition, routed $k$, and router confidence, and which is
read by the representation path and the routers at the same time --- the model
remembers not just what it computed, but who computed it and how sure they were.
On a ladder of 13 models trained on byte-identical data at \TODO{tokens} we
isolate each axis at matched model FLOPs, including the control this
literature keeps forgetting: a \emph{learned} policy against a \emph{random} one
at the same mean budget. \TODO{headline result}
We close with Metis-1.6 Praxis and Logos, two MoRE models pretrained on 1T
tokens, as evidence that the architecture survives contact with a real training run.
\end{abstract}

% =====================================================================
\section{Introduction}

\TODO{Write last. The four beats, in this order:}

\begin{enumerate}
\item Tokens are not equally hard, and every model in production knows this and
does nothing about it. Predicting the second half of \emph{``San''}$\to$\emph{``Francisco''}
gets the same forward pass as resolving a pronoun three paragraphs back.
\item Conditional compute answers this, but every deployed variant turns exactly
one dial --- and ``harder'' is not one-dimensional. A token can need \emph{more
steps} (a multi-step calculation), \emph{more capacity per step} (a rare entity),
or \emph{different capacity} (a switch from parsing to verifying). Depth, width,
and pathway are not substitutes for one another, and a model that can only buy
one of them will buy the wrong one.
\item MoRE learns all three jointly, per token, out of one budget. The claim is
composition and joint training, not any single axis. Adaptive depth has prior
art. Adaptive $k$ has prior art. We say so in Section~\ref{sec:related} rather
than making a reviewer say it for us.
\item We back it with an axis-isolating ladder at matched model FLOPs. The
uncomfortable control --- learned policy versus random policy at the same mean
budget --- is included, and reported whichever way it comes out.
\end{enumerate}

\paragraph{Contributions.}
\begin{itemize}
\item \textbf{MoRE-Core}: joint per-token allocation over recursion depth,
expert width, and expert pathway, with the routing and budget objectives that
make three discrete decisions trainable from task loss (Section~\ref{sec:core}).
\item \textbf{MoRE-RM}: a route-typed recurrent depth memory, native to MoRE in
the strict sense that its key space is the token's own pass and route history,
read by the representation path and the routers simultaneously
(Section~\ref{sec:rm}).
\item A 13-model ablation ladder that isolates each axis at matched model
FLOPs, plus random-policy controls that ask whether the \emph{learned} allocation
does anything or whether any allocation with the same budget would do
(Section~\ref{sec:ablations}).
\item A systems account of dropless packed dynamic compute --- monotonic active
sets, packed layouts, variable-count collectives --- showing the architecture runs
at \TODO{}\% MFU rather than only in principle (Section~\ref{sec:systems}).
\item Two production models, Metis-1.6 Praxis and Logos, pretrained on 1T tokens
(Section~\ref{sec:scale}).
\end{itemize}

% =====================================================================
\section{Related work}
\label{sec:related}

\TODO{This section is adversarial review done in advance. Be maximally generous
to prior art on each individual axis; the novelty claim is the composition, and
it is stronger, not weaker, when the pieces are properly credited. Sober tone
throughout --- no jokes at anyone else's expense.}

\paragraph{Adaptive depth.}
Adaptive Computation Time \CITE{act2016} and PonderNet \CITE{pondernet2021}
established learned halting for recurrent computation. Mixture-of-Depths
\CITE{mod2024} makes per-token layer skipping learnable in transformers.
Mixture-of-Recursions \CITE{mor2025} is the closest prior work on this axis:
per-token adaptive recursion depth over a weight-shared block, trained from
scratch with a recursion router. Ouro \CITE{ouro2025} studies looped models at
scale and reports that recursion improves the \emph{manipulation} of knowledge
rather than the amount stored --- a finding we adopt as a design constraint rather
than dispute (Section~\ref{sec:limitations}).
\textbf{Difference}: these pair adaptive depth with a dense feed-forward block.
Depth is the only dial.

\paragraph{Adaptive width.}
Adaptive per-token expert count is established prior art and we claim no part of
it. ``Harder Tasks Need More Experts'' \CITE{harder2024} states the thesis in its
title; AdaMoE \CITE{adamoe2024} realizes variable $k$ through null experts;
DynaMoE \CITE{dynamoe2026} adds layer-wise adaptive capacity; ProbMoE
\CITE{probmoe2026} formulates routing as a distribution over cardinality-constrained
expert subsets; AnyExperts \CITE{anyexperts2025} allocates on demand in the
multimodal setting. Expert-choice routing \CITE{expertchoice2022} attacks the
same imbalance from the expert side.
\textbf{Difference}: all operate in a single-pass stack. Width is the only dial,
and it is decided once.

\paragraph{Recursion combined with sparse experts.}
This combination became active immediately before this work, and we position
against it directly rather than by omission. LoopMoE \CITE{loopmoe2026} is a
block-recurrent MoE at 3B/9B that resolves weight-sharing symmetry with
IterAdaLN, a modulation signal conditioned jointly on the iteration index and the
token state, evaluated against vanilla MoE at identical total parameters and
per-token FLOPs. The Loopie series \CITE{loopies2026} trains 20B/A2B and 6B/A0.6B
looped MoE models with reasoning post-training. ``Tying the Loop''
\CITE{tyingtheloop2026} compares per-iteration adaptation strategies including
depth-wise LoRA \CITE{relaxedrecursive2024}, low-rank iteration modulation, and
MeSH \CITE{mesh2025}.
\textbf{Difference, stated precisely}: in all of these the loop count is a fixed
hyperparameter and the expert count per iteration is fixed. Neither the depth nor
the width of the computation is a function of the token. Our claim is that making
both adaptive \emph{and jointly budgeted}, with the pathway re-decided from the
evolving state at every pass, is what produces the gain --- and our ladder
includes a fixed-loop fixed-$k$ MoE row (Section~\ref{sec:ablations}, row~5) that
reimplements exactly this design point in our backbone, so the comparison is
internal and matched rather than cross-paper and hopeful.

\paragraph{Residual topology and local memory.}
We use, and do not claim, Hyper-Connections \CITE{hyperconnections2024} in a
manifold-constrained form (mHC), the Mamba-2 hybrid backbone \CITE{mamba2_2024},
$n$-gram conditional memory in the Engram line \CITE{engram}, and aux-loss-free
expert balancing \CITE{deepseekv3}. Section~\ref{sec:integrations} states the one
thing we did modify --- each is conditioned on the pass index --- and claims that
delta and nothing more.

% =====================================================================
\section{MoRE-Core}
\label{sec:core}

\subsection{Setup}

A MoRE model is a weight-shared stack $F$ of $L$ physical layers. Every token
receives pass~1. After pass~$r$, a continuation router decides per token whether
that token halts or enters pass~$r+1$. The active set is monotonic,
$A_{r+1} \subseteq A_r$, and is physically packed at every pass. Depth is capped
at $\Rmax$ passes.

We fix the convention \emph{depth $=$ passes $=$ recursions $+ 1$}, so depth~1 is
the non-looped base case, matching \CITE{ouro2025}.

Each layer's feed-forward sublayer is a mixture of $E$ routed experts plus one
always-active shared expert, computed in a latent space of width $d_\ell < d$.
At every pass, and independently at every layer, a router chooses both how many
routed experts the token gets and which ones.

\subsection{The three decisions}

For token $t$ at pass $r$ with state $h_t^{(r)}$:

\paragraph{Depth.} A continuation router produces
\[
  p_t^{(r)} = \sigma\!\left(g_\theta\!\left(h_t^{(r)},\, m_t^{(r)},\, \Delta_t^{(r)},\, \rho_t^{(r)}\right) + \tau_t^{(r)} - \tfrac{1}{2}\right),
\]
where $m_t^{(r)}$ is retrieved depth memory (Section~\ref{sec:rm}),
$\Delta_t^{(r)} = h_t^{(r)} - h_t^{(r-1)}$ is the state difference,
$\rho_t^{(r)}$ is route history, and $\tau_t^{(r)}$ is a bounded, \emph{detached}
innovation prior $\tanh(\lVert\Delta\rVert / \lVert h \rVert)$. The prior gives
the router a sensible monotonic initialization --- tokens whose state is still
moving get another pass --- while leaving the learned logits free to override it
from task loss.
\TODO{Ablate the prior. A reviewer will reasonably ask whether the depth policy
is learned or whether the prior is quietly doing the work, and they should ask.
This is a cheap eval-time ablation on a trained checkpoint: zero $\tau$ and
re-measure the depth distribution.}

\paragraph{Width.} A $k$-router emits logits over $\{k_{\min},\dots,\kmax\}$. Let
$q_t^{(r)}$ be the softmax at temperature $\gamma$ and
$\hat{k}_t^{(r)} = \sum_j q_{t,j}^{(r)} k_j$ the expected count. The hard counts
are assigned by ranking $\hat{k}$ within the batch under the maximum-entropy
integer marginal whose mean is exactly $k^\star$. The hard count drives the
packed dispatch exactly; easy/low-ranked tokens receive fewer experts and
hard/high-ranked tokens receive more, while neither all-$k^\star$ nor an
all-min/all-max collapse can satisfy the marginal. The routed
contribution is scaled by a straight-through envelope
\[
  \frac{k_t^{(r)} + \hat{k}_t^{(r)} - \operatorname{sg}[\hat{k}_t^{(r)}]}{\operatorname{sg}[k_t^{(r)}]},
\]
identically one in the forward pass, but carrying gradient from the task loss
into the $k$-router. This is what makes width a \emph{learned} decision rather
than a scheduled one.

\paragraph{Exact depth budget.} At every continuation boundary, the model ranks
the active tokens by $p_t^{(r)}$. The number that survive is fixed by the
maximum-entropy integer depth marginal over $\{1,2,3\}$ with mean
$d^\star=2$. Thus the router
learns \emph{which} tokens deserve later passes while the batch budget fixes
only \emph{how many}. Random-depth and random-$k$ controls use the identical
marginals with random rankings, so learned-versus-random rows execute exactly
the same hard compute.

\paragraph{Pathway.} The expert router is re-evaluated at every pass from the
updated state and current route features, so the selected coalition can change
even when $k$ does not. Shared physical weights therefore do not imply repeated
identical computation: successive passes can specialize into stages.
\TODO{This axis has the weakest support in the prior literature and the strongest
story in ours. The expert-transition-matrix figure
(Section~\ref{sec:analysis}) carries it; row~6 of the ladder proves it.}

\subsection{Training three discrete decisions at once}

\[
  \mathcal{L} = \mathcal{L}_{\mathrm{LM}}
  + \lambda_k \big(\bar{\hat{k}} - k^\star\big)^2
  + \lambda_d \big(\bar{\hat{d}} - d^\star\big)^2
  + \lambda_z \mathcal{L}_z
  - \lambda_H \mathcal{H}\!\left[p^{(r)}\right],
\]
with $k^\star$, $d^\star$ the exact hard marginal means; the quadratic terms
calibrate the soft surrogates but do not enforce the hard budget. Expert load balancing
uses the aux-loss-free selection-bias update \CITE{deepseekv3} rather than a
balance loss, so the balancing objective is not fighting the width objective for
control of the same router.
\TODO{Report the calibration coefficients and a sensitivity sweep.}

\paragraph{Why one budget and not two.} The current campaign fixes exact depth
and width marginals separately. That proves learned allocation on both axes but
does \emph{not} by itself prove a joint trade between them. The depth--width
correlation analysis in Section~\ref{sec:analysis} is therefore a falsifier, not
decoration; a true shared-cost controller remains required before the title's
``one budget'' claim can be stated without qualification.

% =====================================================================
\section{MoRE-RM: route-typed recurrent depth memory}
\label{sec:rm}

\subsection{Definition}

At each attention anchor and at the end of every completed pass, each active
token writes a typed memory entry: a learned state representation together with
its \emph{route type} --- pass and anchor identity, the weight-averaged
expert-coalition embedding, routed $k$, expert-router entropy and confidence, and
continuation confidence. At later passes each current residual stream attends
over that token's own valid entries, masked per token and bounded by $\Rmax$.

Retrieval has two simultaneous consumers. The \textbf{representation path} gates
memory back into the residual streams before the mixer, so computation can
recover, compare, verify, or revise earlier reasoning states. The \textbf{routing
path} feeds the depth, width, and pathway heads alongside the current state and
state-difference features. No fixed cosine threshold or hand-written halt rule
decides depth, because none is needed once the routers can see what the earlier
passes actually did.

\subsection{Relation to MeSH and to multi-stream residuals}
\label{sec:mesh}

The nearest published mechanism is MeSH \CITE{mesh2025}, which replaces the
overloaded hidden state of a recursive transformer with an explicit buffer of $B$
slots governed by step-wise read/write routers. Given the surface similarity ---
both add memory to a recursive stack --- the distinction is worth making
precisely rather than gesturing at.

MeSH maintains $B$ slots $\mathbf{m}_b \in \mathbb{R}^{L \times D}$ initialized
with the embeddings in slot~0, and per-iteration routers
$\mathbf{w}^{(t)}_{\mathrm{write}} = \mathrm{softmax}(\mathrm{Linear}^{(t)}_{\mathrm{write}}(\mathbf{h}^{(t)}))$
and likewise for reading, with unique parameters per iteration $t$. Writing is
additive into the slots and reading is a weighted sum over them. Three structural
consequences follow:

\begin{enumerate}
\item \textbf{MeSH is closer to our mHC path than to MoRE-RM.} A set of parallel
per-position state slots, mixed by iteration-conditioned learned weights, is
functionally a multi-stream residual with pass-dependent mixing --- which is what
recursion-aware mHC \CITE{hyperconnections2024} is (Section~\ref{sec:integrations}).
We therefore cite MeSH as convergent with our residual path, and credit it as an
independent derivation of that idea, rather than treating it as prior art for the
memory.
\item \textbf{MeSH slots are untyped and carry no route information.} The buffer
records no expert coalition, no routed $k$, no router entropy, no halting
confidence --- because in a fixed-loop dense-or-MoE-agnostic model none of those
quantities exists. MoRE-RM entries are keyed by exactly those quantities, and the
read is attention over a growing set of typed entries rather than a softmax
mixture over a fixed slot count.
\item \textbf{MeSH memory has no routing consumer.} Its read reconstructs
$\mathbf{h}^{(t+1)}$ and nothing else; recursion depth in MeSH is a fixed
hyperparameter with no halting mechanism, and there is no width or pathway
decision for a memory to inform. MoRE-RM's second consumer --- the routers --- is
not an extra feature bolted on, it is the reason the memory is typed at all.
\end{enumerate}

The honest summary: MeSH solves \emph{undifferentiated computation across
iterations}, and we solve that too, with mHC. MoRE-RM solves a different problem
that only exists once depth and width are decisions --- giving those decisions
evidence about what the earlier passes did and how confident they were. A
single-pass or fixed-loop model cannot adopt MoRE-RM unchanged, which is the
criterion in Section~\ref{sec:integrations} and the reason it is claimed here.

\TODO{Consider one empirical version of this argument: strip the route typing
from MoRE-RM entries (state representation only, no coalition/entropy/confidence)
and re-run. If typed and untyped entries perform identically, the honest
conclusion is that we built a more expensive MeSH, and we should say so. Cheap
enough to be worth knowing.}

% =====================================================================
\section{Integrations we use but do not claim}
\label{sec:integrations}

We apply a single criterion consistently: \emph{a mechanism belongs to the MoRE
contribution if its inputs or key space are defined over passes and routes; if a
single-pass, fixed-budget model could adopt it unchanged, it is an integration.}

\begin{center}
\small
\setlength{\tabcolsep}{4pt}
\begin{tabular}{p{0.42\textwidth}p{0.2\textwidth}p{0.28\textwidth}}
\toprule
Mechanism & Status & Basis \\
\midrule
Depth / width / pathway routing + budget objectives & contribution & definitional \\
Route-typed recurrent depth memory & contribution (MoRE-RM) & keyed by pass and coalition \\
Manifold-constrained Hyper-Connections (mHC) & integration & \CITE{hyperconnections2024} \\
Mamba-2 / attention hybrid backbone & integration & \CITE{mamba2_2024} \\
$n$-gram conditional memory & integration & \CITE{engram} \\
Aux-loss-free expert balancing & integration & \CITE{deepseekv3} \\
\bottomrule
\end{tabular}
\end{center}

Two integrations are modified rather than adopted, and we claim the modification
only. The mHC controller receives a learned pass embedding and pass-specific
biases, so pass~1 through pass~$\Rmax$ use different residual topologies without
duplicating weights. The $n$-gram injection gates are pass-aware, so a retrieved
static vector is not added unconditionally on every pass; the current recurrent
state decides whether to use it. Both deltas exist because the base technique had
no notion of a pass, which is also why they are small.

We note the convergence: LoopMoE's IterAdaLN \CITE{loopmoe2026} and MeSH's
step-wise routers \CITE{mesh2025} are independent solutions to the same
weight-sharing symmetry problem that pass-conditioned mHC addresses. Three
groups reaching for the same fix in the same year is reasonable evidence that the
problem is real.

% =====================================================================
\section{Experiments}
\label{sec:ablations}

\subsection{Design}

All ablation models share one backbone --- the same hybrid mixer stack, mHC
residual path, and $n$-gram memory --- and differ only in the routing axes under
test. Holding the backbone constant means our \emph{baselines are stronger than
their standard counterparts}, which makes the comparison conservative for MoRE
and removes backbone quality as an explanation for any gap we find. We state this
rather than let a reader assume ``Dense'' means a vanilla transformer, because it
does not.

Every model sees the \emph{identical token stream in identical order}, sampled
proportionally from the Metis-1.6 1T mixture including its phase structure, so
data order is not a confound for anyone. Fixed-depth rows use the measured mean
depth of MoRE-Core; fixed-$k$ rows use $k=4$.

\begin{center}
\footnotesize
\setlength{\tabcolsep}{3.5pt}
\begin{tabular}{clccccccr}
\toprule
\# & Model & Recur. & Depth & Width $k$ & Pathway & Mem. & Stored/Active & GF/tok M/E \\
\midrule
1 & Dense, FLOP-matched & -- & 1 & -- & dense & -- & 1.44B / 0.87B & 7.09 / 7.62 \\
2 & Dense, parameter-matched & -- & 1 & -- & dense & -- & 1.81B / 1.24B & 9.31 / 9.85 \\
3 & MoE, $k{=}4$ & -- & 1 & fixed & routed & -- & 1.81B / 0.279B & 3.54 / 4.08 \\
4 & MoE, $k{=}8$ & -- & 1 & fixed & routed & -- & 1.81B / 0.336B & 3.88 / 4.42 \\
5 & Fixed LoopMoE & yes & fixed 2 & fixed & re-routed & -- & 1.81B / 0.279B & 7.09 / 9.45 \\
6 & Loop, pathway frozen & yes & fixed 2 & fixed & frozen & -- & 1.81B / 0.279B & 7.09 / 9.45 \\
7 & MoR + dense FFN & yes & adaptive & -- & dense & -- & 0.85B / 0.278B & 7.08 / 8.15 \\
8 & MoR + fixed-$k$ MoE & yes & adaptive & fixed & re-routed & -- & 1.81B / 0.279B & 7.09 / 8.16 \\
9 & Fixed-depth, adaptive $k$ & yes & fixed 2 & adaptive & re-routed & -- & 1.81B / 0.279B & 7.09 / 9.45 \\
10 & \textbf{MoRE-Core} & yes & adaptive & adaptive & re-routed & -- & 1.81B / 0.279B & 7.09 / 8.16 \\
11 & \textbf{MoRE-RM} & yes & adaptive & adaptive & re-routed & yes & 1.81B / 0.279B & 7.09 / 8.16 \\
\midrule
12 & Random-$k$ control & yes & adaptive & random & re-routed & -- & 1.81B / 0.279B & 7.09 / 8.16 \\
13 & Random-depth control & yes & random & adaptive & re-routed & -- & 1.81B / 0.279B & 7.09 / 8.16 \\
\bottomrule
\end{tabular}
\end{center}

Active parameters are per pass; rows with mean depth 2 execute them roughly
twice. GFLOP/token reports model work / issued hardware work. The scientific
iso-FLOP comparison uses model work; issued work additionally includes the
checkpoint replay required by each measured memory lane. Each row differs from
a neighbour in exactly one factor.

\begin{itemize}
\item \textbf{Rows 5 vs 6} isolate pathway: identical FLOPs, identical depth and
width, differing only in whether the expert coalition is re-decided after pass~1
or frozen. Row~5 reimplements the fixed-loop MoE design point
\CITE{loopmoe2026} in our backbone.
\item \textbf{Rows 8, 9, 10} isolate depth and width against each other at
matched model FLOPs. Row~9 issues more hardware work because its measured
40-APU lane uses layer checkpointing.
\item \textbf{Rows 12, 13} hold the compute budget fixed and randomize the
policy. To our knowledge this control is absent from the adaptive-compute
literature, and it is the experiment most capable of embarrassing us, which is
why it is in the required set rather than the appendix.
\item \textbf{Row 7} is necessarily smaller in stored parameters: a dense-FFN
recursive model cannot hold 1.8B stored at 0.279B active. That asymmetry is not a
flaw in the control, it is precisely what the sparse-expert axis buys, and we
report both numbers rather than picking the flattering one.
\end{itemize}

Rows 1 and 5--13 are iso-model-FLOP by construction. Rows 2, 3, 4 are deliberate
off-frontier reference points. Note that a wider single-pass MoE \emph{cannot} be
made iso-FLOP with a recursive model by raising $k$ alone: going $k{=}4 \to 8$
adds about 0.34 model GFLOP/token because expert GEMMs are a minority of the
block, while a second pass adds roughly 3.55. Row~4 is a reference, not a match, and we say so before
someone assumes we tried and failed.

\subsection{Configuration}

\TODO{Proxy manifest table: 2 physical layers (1 Mamba-2 + 1 attention),
$d_{\mathrm{model}}=4096$, latent $2048$, expert intermediate $1152$, 72 routed + 1
shared, $k \in [1,8]$ with $k^\star=4$, $\Rmax=5$, $d^\star=2$, sequence 4096,
tokenizer 65{,}536. Fill in exact audited counts from the ablation manifest.}

\TODO{Learning rate: report the short per-archetype sweep and which value each
row used. One LR across architectures of different active widths is not fair;
tuning all thirteen is not affordable. State the compromise --- a reviewer who
sees it stated is satisfied; one who suspects it is not.}

\subsection{Scaling}

\TODO{Three sizes $\times$ four archetypes (Dense, MoE, MoRE-Core, MoRE-RM), plus
Praxis and Logos as points four and five on the same curve. The question is
whether the MoRE advantage holds, grows, or shrinks with scale. One size is a
data point; three is a trend; five is a claim.}

\subsection{Results}

\TODO{Everything. Primary figure: quality vs executed FLOPs, all rows, one
frontier plot. Secondary: quality vs tokens. Tertiary: quality vs wall-clock ---
the axis that matters in practice and the one MoRE is most at risk on.}

% =====================================================================
\section{Analysis: does it actually adapt?}
\label{sec:analysis}

Loss curves show MoRE is not worse. This section is the evidence that it does the
thing it says it does.

\begin{itemize}
\item \textbf{Allocation vs difficulty.} Depth and $k$ distributions conditioned
on token surprisal, frequency, and syntactic class.
\item \textbf{Depth--width correlation.} If depth and $k$ move together
perfectly, one axis is redundant and the three-dial framing is oversold. We
report the correlation whatever it says. \TODO{Run this first --- it can change
the paper's central claim, and it is better to find that out in July than in
review.}
\item \textbf{Expert transitions across passes.} Transition matrices between
pass~$r$ and pass~$r{+}1$ coalitions, testing the stage-specialization
hypothesis (parse $\to$ manipulate $\to$ verify).
\item \textbf{Halt calibration.} Reliability diagram for continuation confidence.
\item \textbf{Test-time budget dial.} The same checkpoint under continuation caps
$1..\Rmax$ and forced uniform depths $1..\Rmax$, trained with stochastic depth
budgets so every cap is a supported operating point rather than an
out-of-distribution override.
\item \textbf{Oracle gap.} On a small evaluation subset, exhaustively search the
best per-token depth and report what fraction of the oracle improvement the
learned policy captures.
\end{itemize}

% =====================================================================
\section{Systems: dropless packed dynamic compute}
\label{sec:systems}

Adaptive-compute papers are usually rejected on ``does this actually run fast?''
rather than on quality, and rightly so; a FLOP you cannot schedule is a FLOP you
did not save. This section is a differentiator, not an appendix.

\TODO{Cover: monotonic active sets and per-pass packing; dropless routing with
variable-count collectives; the three quantities that must never be conflated
(training-data token dropout: none; MoE overflow drop rate: zero; learned
continuation halting: intended and load-bearing); FP8 role assignment and the
BF16 parity lane; measured MFU and tokens/second; the cost of dynamic control
flow against a static-shape baseline.}

% =====================================================================
\section{Scale demonstration: Metis-1.6}
\label{sec:scale}

This section is a demonstration, not a controlled comparison, and we would rather
say that in the first sentence than have it inferred from a footnote.

Metis-1.6 Praxis and Metis-1.6 Logos are MoRE models pretrained on the same
immutable 1T-token release with identical phase boundaries, tokenizer, and
curriculum, differing only in scale. \TODO{parameter and active-parameter audits}

\paragraph{Comparison to published models.} Any comparison across laboratories
differs in data, token budget, and post-training, so we report total parameters,
active parameters per token, and training tokens beside every score, and we plot
quality against active FLOPs per token rather than presenting a ranked table.

We state one prediction in advance: at a 1T-token budget, MoRE models should be
\emph{competitive on reasoning and manipulation benchmarks and behind on
closed-book knowledge benchmarks} relative to models trained on an order of
magnitude more tokens. This follows directly from the depth-versus-knowledge
finding in \CITE{ouro2025} --- recursion buys manipulation, tokens buy knowledge
--- and we would rather register it as a prediction of the architecture now than
offer it as an explanation later.

% =====================================================================
\section{Claims, evidence, and falsifiers}

\begin{center}
\begin{tabular}{p{0.29\textwidth}p{0.3\textwidth}p{0.33\textwidth}}
\toprule
Claim & Evidence & Falsified if \\
\midrule
Depth adaptivity beats a matched fixed depth & rows 9 vs 10 at iso-model-FLOP & row 9 $\geq$ row 10 \\
Width adaptivity beats a matched fixed $k$ & rows 8 vs 10 at iso-model-FLOP & row 8 $\geq$ row 10 \\
Per-pass re-routing matters & rows 5 vs 6 & no gap at iso-FLOP \\
The learned policy beats a random one at the same budget & rows 12, 13 vs 10 & no gap \\
Depth and width are distinct axes & correlation analysis & near-perfect correlation \\
Later passes do useful work & per-pass quality gain & gain vanishes after pass 2 \\
Route memory helps & row 11 vs 10 & no gap, or later experts starve \\
Route \emph{typing} is what helps & typed vs untyped entries & no gap (we built a costly MeSH) \\
Stage specialization is real & transition matrices & coalitions are pass-invariant \\
\bottomrule
\end{tabular}
\end{center}

These were written before the runs, and we report against them either way.

% =====================================================================
\section{Limitations}
\label{sec:limitations}

\begin{itemize}
\item \textbf{Ablation scale.} \TODO{state tokens/parameter per row}. Conclusions
about relative architecture quality at 1.5B do not automatically transfer to
frontier scale. The Praxis and Logos points partially address this and do not
close it, and anyone claiming otherwise about their own ablations is guessing too.
\item \textbf{Recursion does not buy knowledge} \CITE{ouro2025}. MoRE targets
manipulation. Closed-book factual benchmarks are not where this architecture is
expected to win, and we do not report them as though they were.
\item \textbf{Wall-clock is not FLOPs.} Dynamic routing costs launch overhead,
compaction, and variable collectives. We report wall-clock next to FLOPs even
where the comparison is less flattering.
\item \textbf{Budget coefficients are choices.} Two quadratic penalties on a
shared budget is the most tunable part of the design. \TODO{sensitivity sweep}
\item \textbf{The ladder shares a non-standard backbone.} Conclusions are about
the routing axes on that backbone, and we have not shown they transfer to a
vanilla transformer stack.
\item \textbf{Three dials mean three ways to be wrong.} Every axis we added is
another policy that can fail to learn, and the analysis section exists because we
did not want to find that out from a reader.
\end{itemize}

% =====================================================================
\section{Conclusion}
\TODO{}

\bibliographystyle{plain}
\bibliography{references}

\end{document}
